\section{Introducción}
\label{sec:introduccion}

\subsection{Antecedentes}
\label{sec:antecedentes}

La clasificación espectral es uno de los pilares de la astrofísica estelar
moderna. El sistema de Morgan--Keenan (MK) \citep{morgan1943} ordena a las estrellas en una
secuencia de tipos espectrales (O, B, A, F, G, K, M), de las más calientes y
azules a las más frías y rojas, complementada por una clase de luminosidad
(I--V) que distingue gigantes de estrellas de secuencia principal. Cada tipo se
correlaciona con la temperatura efectiva (\teff) de la fotosfera y, junto con la
clase de luminosidad, codifica de forma compacta los parámetros físicos
fundamentales de un objeto: temperatura, gravedad superficial (\logg) y, de
manera indirecta, su estado evolutivo. A pesar de tener más de un siglo de
antigüedad, el esquema MK sigue siendo el lenguaje de referencia con el que se
comunican los resultados de la espectroscopía estelar.

El problema de asignar un tipo MK a partir de observaciones reales dista de
ser trivial. Los espectros y fotometría reales presentan ruido instrumental,
\emph{enrojecimiento} por extinción interestelar, y contaminación por
binariedad o por fuentes problemáticas, factores que confunden a los métodos
tradicionales de plantillas y de ajuste de continuo. La estimación fotométrica
de la extinction $E(BP-RP)$ \verb|ebpminrp_gspphot|, por ejemplo, es
propensa a valores ausentes (\texttt{NaN}) y arrastra incertidumbres
significativas, especialmente para estrellas frías o de baja gravedad
superficial. Cualquier sistema de clasificación que aspire a operar sobre datos
reales debe manejar de manera explícita estas anomalías (
paralajes negativas, campos químicos vacíos y banderas de calidad), en lugar
de suponer entradas limpias.

\subsubsection{Gaia DR3}

La tercera entrega de datos de la misión \emph{Gaia} de la Agencia Espacial
Europea (Gaia DR3) marcó un cambio de escala en la astronomía estelar al
publicar astrometría, fotometría y parámetros astrofísicos para más de mil
millones de fuentes. Para una fracción importante de ellas, los \emph{pipelines}
de procesamiento de Gaia (GSP-Phot, GSP-Spec, ESP-ELS) entregan no sólo
posiciones y movimientos propios, sino también estimaciones de \teff, \logg,
metalicidad global $[M/H]$, abundancia de hierro y anchos equivalentes de
líneas como H$\alpha$. Este volumen de datos hace inviable la clasificación
visual estrella por estrella que históricamente realizaban los expertos, y crea
la necesidad de sistemas automáticos que sean, a la vez: precisos,
reproducibles y \emph{explicables}.

Es justo la última necesidad, la de explicabilidad, uno de los principales motores
que impulsan la propuesta de este proyecto. Dentro del campo de la astronomía, los
expertos no necesariamente dominan los detalles técnicos y teóricos de cómo funcionan 
los modelos de Deep Learning. Es por ello que la capacidad de auditar el razonamiento
fue fundamental para la construcción de STELLAR, y de cualquier otro sistema que se 
propuesiera para resolver este problema. En ese sentido, los modelos de redes neuronales,
los modelos de lenguaje a gran escala, y algunas otras arquitecturas de Inteligencia 
Artificial, pueden llegar a ser denominados como "cajas negras", no porque no se pueda
entender su funcionamiento, sino porque la toma de decisiones no es tan transparente
para un investigador que no esté familiarizado con su funcionamiento. 

El corpus de trabajo de este proyecto se construye precisamente sobre este
acervo: un conjunto curado de \num{498} estrellas de Gaia DR3 \citep{gaiadr3}, con su
metadata tabular escalar (astrometría, fotometría, química y física derivada).
El catálogo respeta de forma estricta las convenciones del archivo, incluyendo
la convención de signo del pipeline ESP-ELS para H$\alpha$, en la que los
valores positivos indican absorción y los negativos emisión real.

\subsubsection{Modelos de lenguaje especializados en astronomía}

Paralelamente, los \textit{Large Language Models} (LLM), conocidos en español como
modelos de lenguaje a gran escala, han demostrado capacidad
de razonamiento sobre dominios técnicos cuando se les ajusta (\emph{fine-tuning})
con corpus especializados. En astronomía, modelos como \AstroSage{} \citep{astrosage} (un LLM de
\num{8} mil millones de parámetros afinado sobre literatura astronómica),
abren la posibilidad de emitir diagnósticos en lenguaje natural que no sólo
asignan una etiqueta, sino que articulan el razonamiento astrofísico que la
sustenta. Esta capacidad de explicación es especialmente valiosa en un dominio
donde la confianza del experto depende de poder auditar el porqué de cada
clasificación.

Sin embargo, emplear un LLM como clasificador directo sobre datos crudos tiene
riesgos bien documentados: alucinaciones, sensibilidad al formato de entrada y
sesgos \emph{a priori} hacia las clases más frecuentes en su corpus de
entrenamiento. En las observaciones preliminares de este trabajo se identificó
precisamente un sesgo fuerte del modelo hacia el tipo G (estrellas similares al
Sol), que distorsiona la clasificación cuando el modelo opera sin restricciones.

\subsubsection{Hacia una arquitectura híbrida HC + SC}

Estos antecedentes motivan la propuesta central de este reporte: una
arquitectura híbrida de dos capas que combina \emph{Hard Computing} (HC) y
\emph{Soft Computing} (SC). La capa HC aplica algoritmos deterministas basados
en física clásica y estadística para extraer parámetros limpios y generar un
conjunto de \emph{banderas lógicas} (una ``predigestión'' de la evidencia
observacional: paralaje confiable, candidato a binaria, baja metalicidad,
etc.). La capa SC entrega ese contrato estructurado al modelo \AstroSage{}, que
produce la clasificación MK final y una descripción en lenguaje natural. De
este modo, el anclaje determinista del HC acota el espacio de respuestas del LLM
y corrige su sesgo \emph{a priori}, mientras que el SC aporta la capacidad de
razonamiento y explicación. El sistema resultante, denominado \STELLAR{}
(\emph{Spectral Type Estimation via Language Learning and Astronomical
Reasoning}), es el objeto de estudio de este reporte.



\subsection{Planteamiento del Problema}
\label{sec:problema}

El problema que aborda este trabajo puede enunciarse de la siguiente manera:
\emph{dado el vector de metadata tabular de una estrella de Gaia DR3
(astrometría, fotometría, química y parámetros físicos derivados),
producir de forma automática (i)~una clasificación espectral MK , i.e. tipo,
subtipo y clase de luminosidad; (ii)~una caracterización de su población
galáctica, y (iii)~una justificación astrofísica en lenguaje natural, de manera
precisa, reproducible y auditable.}

Resolver este problema con un único paradigma de cómputo conduce a soluciones
insatisfactorias en ambos extremos del espectro Hard Computing--Soft Computing:

\begin{itemize}
  \item \textbf{Enfoque puramente determinista (sólo HC).} Codificar todo el
  conocimiento experto de la clasificación MK en reglas y umbrales rígidos
  resulta frágil ante datos reales. Las fronteras entre tipos espectrales
  adyacentes son difusas (la separación B$\leftrightarrow$A ocurre cerca de
  $T_\mathrm{eff}\approx\SI{10000}{\kelvin}$, donde un corte duro genera errores
  sistemáticos), y un sistema de reglas no produce, por construcción, la
  explicación en lenguaje natural que demanda el experto. Además, debe absorber
  toda la suciedad del dato: valores \texttt{NaN}, paralajes negativas, banderas
  de calidad o campos químicos vacíos.

  \item \textbf{Enfoque puramente generativo (sólo SC).} Entregar los datos
  crudos directamente a un modelo de lenguaje hereda los riesgos conocidos de
  los LLM: alucinaciones, sensibilidad al formato de la entrada y, de manera
  crítica para este dominio, un sesgo \emph{a priori} hacia las clases más
  frecuentes en su corpus de entrenamiento. En las pruebas preliminares se
  observó un prior fuerte hacia el tipo G (estrellas solares), que sesga las
  predicciones y degrada la exactitud sobre los extremos calientes y fríos de
  la secuencia. El LLM tampoco garantiza, por sí solo, el respeto de
  convenciones técnicas como el signo de H$\alpha$ en el pipeline ESP-ELS.
\end{itemize}

De esta tensión se desprende el problema metodológico central: \emph{¿cómo
combinar el rigor determinista del Hard Computing con la capacidad de razonamiento y
explicación del Soft Computing, de modo que el primero ancle y acote al segundo 
sin anular su flexibilidad?} En concreto, el trabajo busca responder:

\begin{enumerate}
  \item \textbf{Anclaje.} ¿Puede una capa HC determinista, que predigiere la
  evidencia observacional en un contrato estructurado con banderas lógicas,
  corregir el sesgo \emph{a priori} del LLM y elevar la exactitud de la
  clasificación MK?
  \item \textbf{Calidad de la clasificación.} ¿Qué nivel de exactitud, acuerdo
  (Cohen $\kappa$) y cercanía (\emph{near-miss}) alcanza el sistema híbrido
  frente a las etiquetas de referencia de SIMBAD, y cómo se distribuyen sus
  errores entre clases espectrales?
  \item \textbf{Consistencia física.} ¿Son los parámetros inferidos (\teff,
  \logg) consistentes con catálogos de alta resolución de la literatura
  (PASTEL)?
  \item \textbf{Calidad de la explicación.} ¿Son las descripciones en lenguaje
  natural coherentes con la clasificación emitida y semánticamente fieles a
  descripciones de referencia?
\end{enumerate}

La validación se realiza sobre el corpus completo de \num{498} estrellas de
Gaia DR3, evaluando de forma iterativa siete versiones del \emph{system prompt}
de la capa SC (V1--V7) para aislar el efecto de cada decisión de diseño sobre
las métricas anteriores.

\subsection{Justificación}
\label{sec:justificacion}

La pertinencia de \STELLAR{} se sostiene sobre cuatro ejes: la escala del
problema astronómico, la necesidad de explicabilidad, el valor metodológico del sistema 
híbrido (HC+SC), y los resultados empíricos obtenidos.

\paragraph{Escala y automatización.} Los sondeos masivos como Gaia han
desplazado el cuello de botella de la astronomía estelar desde la obtención de
datos hacia su interpretación. Con parámetros astrofísicos publicados para
cientos de millones de fuentes, la clasificación visual experta, estrella por
estrella, es claramente inviable. Un sistema automático que opere directamente sobre 
la metadata tabular de Gaia DR3, y que tolere de forma robusta sus anomalías
(\texttt{NaN}, paralajes negativas, banderas de calidad, campos químicos
vacíos), responde a una necesidad concreta y creciente de la disciplina.

\paragraph{Explicabilidad como requisito.} En astrofísica la
confianza del especialista depende de poder auditar el razonamiento detrás de
cada etiqueta. Un clasificador de ``caja negra'' que sólo emite ``tipo K'' aporta
poco frente a uno que justifica la decisión en términos de temperatura,
gravedad, metalicidad y posición en el diagrama HR. Al situar un LLM
especializado (\AstroSage{}) en la capa de decisión, \STELLAR{} produce, junto
a la clasificación MK, una descripción en lenguaje natural auditable. Esto
convierte al sistema en una herramienta de apoyo a la decisión y no en un mero
predictor de etiquetas.

\paragraph{Aporte metodológico de la hibridación.} La contribución central de
este trabajo no es un nuevo algoritmo aislado, sino una \emph{arquitectura} que
articula dos paradigmas habitualmente disjuntos. La capa HC determinista
predigiere la evidencia en un contrato estructurado y un conjunto de banderas
lógicas que \emph{anclan} al modelo de lenguaje, acotando su espacio de
respuestas y corrigiendo su sesgo \emph{a priori} hacia el tipo G. El patrón
resultante (HC como ancla verificable, SC como motor de razonamiento y
explicación), es transferible a otros problemas científicos donde coexisten
reglas físicas duras y la necesidad de interpretación flexible. Lo anterior constituye en
sí mismo un resultado reutilizable más allá del caso estelar y astronómico.

\paragraph{Evidencia empírica.} La justificación no es sólo conceptual: el
sistema alcanza una exactitud de clasificación MK de \num{0.7951} con un acuerdo
Cohen $\kappa = \num{0.7529}$ y una exactitud \emph{near-miss} ($d\le 1$) de
\num{0.998} sobre las \num{498} estrellas del corpus, manteniendo una desviación
media de \teff{} frente a PASTEL del orden de \SI{200}{\kelvin}.\footnote{Cifras
de la validación V7; el detalle por versión y por clase se presenta en la
Sección~\ref{sec:resultados}.} Estos números confirman que el anclaje
determinista eleva la calidad de la clasificación respecto al uso directo del
LLM, validando la hipótesis de diseño.

\paragraph{Reproducibilidad.} Finalmente, el proyecto se ciñe a una práctica de
reproducibilidad explícita: corpus de tamaño fijo y documentado ($n=498$),
convenciones de datos declaradas (signo de H$\alpha$ ESP-ELS, tratamiento de
\texttt{NaN} y del artefacto $\mathrm{fe\_h}=0.0$), versionado de los
\emph{prompts} (V1--V7), métricas reportadas con intervalos de confianza
\emph{bootstrap} e infraestructura de cómputo especificada. Esto permite que los
resultados sean verificables y que el sistema pueda ser extendido a muestras
mayores de Gaia DR3/DR4.
